\documentclass[11pt,a4paper]{article}

% --- Packages ---
\usepackage[margin=2.5cm]{geometry}
\usepackage{booktabs}
\usepackage{tabularx}
\usepackage{array}
\usepackage{graphicx}
\usepackage{hyperref}
\usepackage{amsmath}
\usepackage{xcolor}
\usepackage{parskip}
\usepackage{enumitem}
\usepackage{titlesec}
\usepackage{fancyhdr}
\usepackage{microtype}

% --- Page style ---
\pagestyle{fancy}
\fancyhf{}
\rhead{Master Thesis Progress Report}
\lhead{Ali Azizpourian}
\cfoot{\thepage}

% --- Title ---
\title{\textbf{Progress Report}\\[0.4em]
       \large Key-Value Pair Extraction from Document Images\\[0.2em]
       \normalsize Master Thesis --- Supervisor Update}
\author{Ali Azizpourian}
\date{April 2026}

\begin{document}

\maketitle
\thispagestyle{fancy}

% --- Overview ---
\section*{Overview}

This report summarises work completed on the master thesis project:
\textit{Key-Value Pair Extraction from Document Images using LayoutLMv3 with a
Biaffine Linker}.
Stages 1 through 3 are complete and have produced the full prepared dataset
ready for model training.
Stage 4 (model fine-tuning) is the next planned phase and is described in
Section~\ref{sec:stage4}.

The pipeline processes raw document images from the IBM KVP10k benchmark
dataset and produces, step by step, structured key-value predictions at the
word level.

% --- Stage 1 ---
\section{Stage 1 --- Dataset Acquisition}

\subsection*{Objective}
Obtain the IBM KVP10k benchmark and establish a reproducible data loading
pipeline.

\subsection*{Dataset}
KVP10k~\cite{kvp10k} is a large-scale document understanding benchmark
published by IBM Research. It contains scanned and born-digital business
documents (invoices, receipts, purchase orders, forms) with ground-truth
key-value pair annotations at the bounding-box level.

\begin{table}[h!]
  \centering
  \caption{IBM KVP10k dataset statistics}
  \label{tab:dataset}
  \begin{tabular}{lrr}
    \toprule
    \textbf{Split} & \textbf{Documents} & \textbf{Annotated KVP entries} \\
    \midrule
    Train          & 48{,}280           & $\sim$640{,}000 \\
    Test           & 5{,}255            & $\sim$70{,}000 \\
    \midrule
    Total          & 53{,}535           & $\sim$710{,}000 \\
    \bottomrule
  \end{tabular}
\end{table}

\subsection*{KVP Entry Types}
Each annotated entry carries one of three semantic types:
\begin{itemize}[noitemsep]
  \item \texttt{kvp} --- a complete key-value pair (both key and value present)
  \item \texttt{unvalued} --- a key with no associated value
  \item \texttt{unkeyed} --- a value with no associated key
\end{itemize}

\subsection*{Output}
Raw JSON annotation files and document images downloaded and cached locally
on the HPC cluster.

% --- Stage 2 ---
\section{Stage 2 --- Document Preprocessing}

\subsection*{Objective}
Convert raw document PDFs to standardised page images and extract layout
metadata required for spatial modelling.

\subsection*{Steps}
\begin{enumerate}[noitemsep]
  \item PDF pages rendered to PNG images at a fixed resolution.
  \item Page dimensions (\texttt{image\_width}, \texttt{image\_height} in
        pixels) recorded per document.
  \item Bounding-box coordinates kept in original pixel space for
        traceability; normalisation to $[0, 1000]$ is performed later
        (Stage 3) as required by LayoutLMv3.
\end{enumerate}

\subsection*{Output}
PNG images and per-page metadata files stored under \texttt{data/raw/}.

% --- Stage 3 ---
\section{Stage 3 --- Data Preparation (LMDX Encoding)}

\subsection*{Objective}
Transform raw annotations into a format directly consumable by the
LayoutLMv3 processor, producing one self-contained JSON file per document.

\subsection*{LMDX Text Format}
Each word is encoded on a single line as:
\begin{center}
  \texttt{word\_text\;x1|y1|x2|y2}
\end{center}
where coordinates are in pixel space. The full page is serialised as a
newline-delimited sequence of such lines stored in the field
\texttt{lmdx\_text}.

\subsection*{Prepared JSON Schema}

\begin{table}[h!]
  \centering
  \caption{Prepared JSON file schema}
  \label{tab:json}
  \begin{tabularx}{\textwidth}{llX}
    \toprule
    \textbf{Field} & \textbf{Type} & \textbf{Description} \\
    \midrule
    \texttt{hash\_name}   & string & Unique SHA-256 document identifier \\
    \texttt{lmdx\_text}   & string & Words with pixel-space bounding boxes
                                      (one word per line,
                                      format \texttt{word x1|y1|x2|y2}) \\
    \texttt{image\_width} & int    & Page width in pixels \\
    \texttt{image\_height}& int    & Page height in pixels \\
    \texttt{gt\_kvps}     & object & Ground-truth KVPs; contains
                                      \texttt{kvps\_list} --- a list of
                                      entries each with \texttt{type},
                                      \texttt{key.text}, \texttt{key.bbox},
                                      \texttt{value.text},
                                      \texttt{value.bbox} \\
    \bottomrule
  \end{tabularx}
\end{table}

\subsection*{Prepared Dataset Statistics}

\begin{table}[h!]
  \centering
  \caption{Prepared dataset split statistics (training set analysed across all 5{,}389 files)}
  \label{tab:prepared}
  \begin{tabular}{lrrrr}
    \toprule
    \textbf{Split} & \textbf{JSON files} & \textbf{Total KVP entries}
    & \textbf{Valid \texttt{kvp} pairs} & \textbf{Partial entries} \\
    \midrule
    Train & 5{,}389 & 81{,}817 & 35{,}363\;(43.2\,\%) & 46{,}454\;(56.8\,\%) \\
    Test  & $\sim$600 & --- & --- & --- \\
    \bottomrule
  \end{tabular}
\end{table}

\noindent
A key data characteristic identified during preparation: only $43.2\,\%$
of annotated entries are complete \texttt{kvp} pairs with non-empty text
and valid bounding boxes. The remaining $56.8\,\%$ are either
\texttt{unvalued} keys (17{,}564 entries) or \texttt{unkeyed} values
(24{,}593 entries) --- partial annotations that are valid for entity
labelling but must not generate key-value link targets. This distinction
is handled explicitly in the data preparation step for Stage 4.

\subsection*{Output}
\begin{itemize}[noitemsep]
  \item \texttt{data/prepared/train/*.json} --- 5{,}389 files
  \item \texttt{data/prepared/test/*.json} --- held-out test set
\end{itemize}

% --- Stage 4 ---
\section{Stage 4 --- Model Fine-Tuning (Planned)}
\label{sec:stage4}

\subsection*{Objective}
Fine-tune \textbf{LayoutLMv3-base}~\cite{layoutlmv3} (Microsoft, 125M
parameters) on the prepared KVP10k data for end-to-end key-value pair
extraction using a two-component architecture.

\subsection*{Model Architecture}
\begin{enumerate}[noitemsep]
  \item \textbf{LayoutLMv3 Encoder} --- joint text, layout (bounding box),
        and visual (patch pixel) encoding; output: $[B,\,512,\,768]$
        hidden states.
  \item \textbf{Entity Classifier} --- token-level linear head classifying
        each token as \textit{Other} (0), \textit{Key} (1), or
        \textit{Value} (2) via cross-entropy loss.
  \item \textbf{Biaffine Linker} --- pairs key tokens to value tokens using
        scaled dot-product semantic scoring combined with an 8-feature
        spatial MLP; trained with binary cross-entropy on the
        key$\times$value adjacency matrix.
\end{enumerate}

\subsection*{Training Configuration}

\begin{table}[h!]
  \centering
  \caption{Planned Stage 4 training hyperparameters}
  \label{tab:hparams}
  \begin{tabular}{ll}
    \toprule
    \textbf{Hyperparameter}          & \textbf{Value} \\
    \midrule
    Base model                       & \texttt{microsoft/layoutlmv3-base} \\
    Max sequence length              & 512 tokens \\
    Learning rate                    & $2 \times 10^{-5}$ \\
    Batch size                       & 1 (gradient accumulation: 8 steps) \\
    Effective batch size             & 8 \\
    Optimiser                        & AdamW (weight decay $= 0.01$) \\
    LR schedule                      & Linear warmup (500 steps) + decay \\
    Max epochs                       & 10 \\
    Early stopping patience          & 3 epochs \\
    Linker loss weight $\lambda$     & Ablation: $\{0.5,\;1.0,\;2.0\}$ \\
    Hardware                         & NVIDIA A100 40\,GB (SLURM cluster) \\
    Mixed precision                  & fp32 \\
    \bottomrule
  \end{tabular}
\end{table}

\subsection*{Two-Stage Training Strategy}
\begin{itemize}[noitemsep]
  \item \textbf{Stage 4a} --- Train encoder + entity classifier only
        (linker disabled). Produces a pretrained encoder checkpoint
        used to warm-start Stage 4b.
  \item \textbf{Stage 4b} --- Load Stage 4a checkpoint; train full model
        including the biaffine linker. The linker loss weight $\lambda$
        is swept across $\{0.5, 1.0, 2.0\}$ to study its effect on
        pair-linking accuracy.
\end{itemize}

\subsection*{Evaluation Metrics}
\begin{itemize}[noitemsep]
  \item \textbf{Entity F1} --- precision and recall on Key/Value token
        classification.
  \item \textbf{KVP F1} --- end-to-end metric: a predicted pair is correct
        if the key text matches (NED $< 0.2$) and bounding boxes overlap
        sufficiently (IoU $> 0.3$).
  \item \textbf{Linker precision/recall} --- standalone linker evaluation
        on correctly classified key-value candidates.
\end{itemize}

\subsection*{Anticipated Challenges}
\begin{itemize}[noitemsep]
  \item \textbf{Data preprocessing complexity} --- the distinction between
        \texttt{kvp}, \texttt{unvalued}, and \texttt{unkeyed} entries
        requires careful label generation: entity labels are assigned for
        all three types, but link labels only for complete \texttt{kvp}
        entries with non-empty text and non-degenerate bounding boxes.
  \item \textbf{GPU memory constraints} --- the 512-token limit and the
        $512 \times 512$ link label matrix impose strict batch-size
        restrictions on the 40\,GB A100; chunked linker scoring and
        a hard candidate cap (max 128 keys / 128 values per page) are
        applied to remain within the memory budget.
  \item \textbf{Training stability} --- the biaffine linker introduces
        additional gradient pathways that require careful learning-rate
        tuning and logit clamping to prevent numerical instability.
\end{itemize}

% --- Bibliography ---
\bibliographystyle{plain}
\begin{thebibliography}{9}

\bibitem{kvp10k}
IBM Research,
\textit{KVP10k: A Comprehensive Dataset for Key-Value Pair Extraction
from Business Documents}, 2023.\\
\url{https://huggingface.co/datasets/ibm/KVP10k}

\bibitem{layoutlmv3}
Y.~Huang et al.,
\textit{LayoutLMv3: Pre-Training for Document AI with Unified Text and
Image Masking},
ACM Multimedia, 2022.

\end{thebibliography}

\end{document}
